% Общая преамбула конспектов (ручной файл; основа -- hse-uberhandbook/common/prelude.tex).
\documentclass[12pt,a4paper]{report}

\usepackage{fontspec}
\usepackage{polyglossia}
\setmainlanguage{russian}
\setotherlanguage{english}

\setmainfont{IBM Plex Serif}
\setsansfont{IBM Plex Sans}
\setmonofont{IBM Plex Mono}

\usepackage[a4paper,margin=12mm,headheight=15pt,headsep=6pt,footskip=8mm]{geometry}
\usepackage{microtype}
\usepackage{setspace}
\setstretch{1.08}
\frenchspacing
\setlength{\parindent}{0pt}
\setlength{\parskip}{4pt plus 1pt}

\usepackage[shortlabels]{enumitem}
\setlist{noitemsep, topsep=3pt, leftmargin=20pt}

\usepackage{graphicx}
\usepackage{fancyvrb}   % блоки кода из #+begin_src
\usepackage{caption}

% Заголовки глав и разделов без выключки по ширине: длинные названия с
% «Армихо-Вульфа» иначе вылезают на поля.
\usepackage[raggedright]{titlesec}

% Колонтитулы: глава слева, раздел справа, номер страницы внизу; каждый
% обрезается до половины строки, чтобы длинные названия не наезжали.
\usepackage{fancyhdr}
\usepackage[fit]{truncate}
\pagestyle{fancy}
\fancyhf{}
\renewcommand{\chaptermark}[1]{\markboth{\thechapter.\ #1}{}}
\renewcommand{\sectionmark}[1]{\markright{\thesection.\ #1}}
\fancyhead[L]{\small\nouppercase{\truncate{0.49\textwidth}{\leftmark}}}
\fancyhead[R]{\small\nouppercase{\truncate{0.49\textwidth}{\hfill\rightmark}}}
\fancyfoot[C]{\thepage}
\renewcommand{\headrulewidth}{0.3pt}
\fancypagestyle{plain}{\fancyhf{}\fancyfoot[C]{\thepage}\renewcommand{\headrulewidth}{0pt}}

\usepackage{hyperref}
\hypersetup{
  colorlinks=true,
  linkcolor=blue!50!black,
  urlcolor=blue!50!black,
  citecolor=blue!50!black,
  bookmarksnumbered=true,
  bookmarksopen=true,
}

\usepackage{amsmath}
\allowdisplaybreaks

\usepackage[nameinlink,noabbrev]{cleveref}

\crefname{theorem}{теорема}{теоремы}
\Crefname{theorem}{Теорема}{Теоремы}

\crefname{lemma}{лемма}{леммы}
\Crefname{lemma}{Лемма}{Леммы}

\crefname{corollary}{следствие}{следствия}
\Crefname{corollary}{Следствие}{Следствия}

\crefname{definition}{определение}{определения}
\Crefname{definition}{Определение}{Определения}

\crefname{proposition}{утверждение}{утверждения}
\Crefname{proposition}{Утверждение}{Утверждения}

\crefname{method}{метод}{методы}
\Crefname{method}{Метод}{Методы}

\crefname{concept}{понятие}{понятия}
\Crefname{concept}{Понятие}{Понятия}

\crefname{example}{пример}{примеры}
\Crefname{example}{Пример}{Примеры}

\crefname{remark}{замечание}{замечания}
\Crefname{remark}{Замечание}{Замечания}

\crefname{equation}{}{}
\Crefname{equation}{}{}

\crefname{chapter}{глава}{главы}
\Crefname{chapter}{Глава}{Главы}

\crefname{section}{раздел}{разделы}
\Crefname{section}{Раздел}{Разделы}

\crefname{subsection}{подраздел}{подразделы}
\Crefname{subsection}{Подраздел}{Подразделы}
